% ===== PRÉAMBULE COMMUN — À COPIER TEL QUEL EN TÊTE DE CHAQUE FICHIER .tex =====
% Ne modifier QUE :
%   - les deux lignes \fancyhead[L] et \fancyhead[R]
%   - le bloc titre \begin{center}...\end{center} juste après \begin{document}
\documentclass[11pt,a4paper]{article}
\usepackage[utf8]{inputenc}
\usepackage[T1]{fontenc}
\usepackage{lmodern}
\usepackage[french]{babel}
\usepackage{amsmath,amssymb,mathtools}
\usepackage{icomma}
\usepackage{xcolor}
\usepackage{tikz}
\usepackage{pgfplots}
\usepackage[most]{tcolorbox}
\usepackage{enumitem}
\usepackage{array,booktabs,colortbl}
\usepackage[a4paper,margin=2.2cm,headheight=26pt,headsep=10pt]{geometry}
\usepackage{fancyhdr}
\usepackage[hidelinks]{hyperref}
\usetikzlibrary{arrows.meta,calc,angles,quotes,positioning,decorations.markings,intersections,backgrounds}
\pgfplotsset{compat=1.18}

\definecolor{primary}{RGB}{226,74,88}
\definecolor{primarydark}{RGB}{150,28,42}
\definecolor{methcol}{RGB}{40,90,150}
\definecolor{excol}{RGB}{0,135,90}
\definecolor{defcol}{RGB}{0,114,178}
\definecolor{attncol}{RGB}{200,40,55}
\definecolor{thmcol}{RGB}{0,140,120}
\definecolor{courbe}{RGB}{32,110,200}
\definecolor{courbeder}{RGB}{210,70,40}
\definecolor{tang}{RGB}{0,150,90}
\definecolor{grille}{RGB}{200,205,215}
\definecolor{plus}{RGB}{200,60,60}
\definecolor{moins}{RGB}{30,90,200}

\tcbset{boxbase/.style={enhanced,breakable,boxrule=0.9pt,arc=2.5pt,
  left=3mm,right=3mm,top=2.3mm,bottom=2.3mm,fonttitle=\bfseries,coltitle=white}}
\newtcolorbox{cle}[1][]{boxbase,colback=defcol!6,colframe=defcol,
  title={\textsc{À retenir}\ifx\relax#1\relax\relax\else\textmd{ : \itshape #1}\fi}}
\newtcolorbox{methode}[1][]{boxbase,colback=methcol!7,colframe=methcol,sharp corners,
  title={\textsc{Méthode}\ifx\relax#1\relax\relax\else\textmd{ --- \itshape #1}\fi}}
\newtcolorbox{exemplebox}[1][]{boxbase,colback=excol!5,colframe=excol,
  title={\textsc{Exemple}\ifx\relax#1\relax\relax\else\textmd{ : \itshape #1}\fi}}
\newtcolorbox{piege}[1][]{boxbase,colback=attncol!6,colframe=attncol,
  title={\textsc{Attention}\ifx\relax#1\relax\relax\else\textmd{ : \itshape #1}\fi}}
\newtcolorbox{exo}[1][]{boxbase,colback=primary!4,colframe=primary,
  title={\textsc{Exercice}\ifx\relax#1\relax\relax\else\textmd{~#1}\fi}}
\newtcolorbox{corr}[1][]{boxbase,colback=thmcol!5,colframe=thmcol,
  title={\textsc{Corrigé}\ifx\relax#1\relax\relax\else\textmd{~#1}\fi}}

\newcommand{\R}{\mathbb{R}}
\newcommand{\N}{\mathbb{N}}
\newcommand{\ds}{\displaystyle}
\newcommand{\tg}{\operatorname{tg}}
\newcommand{\cotg}{\operatorname{cotg}}
\newcommand{\arctg}{\operatorname{arctg}}
\newcommand{\arccotg}{\operatorname{arccotg}}
\newcommand{\limh}{\lim_{h\to 0}}

\pagestyle{fancy}\fancyhf{}
\fancyhead[L]{\small\textcolor{primarydark}{\textbf{Thème 3} — Fonctions composées}}
\fancyhead[R]{\small\textcolor{primarydark}{\textsc{Corrigé — Facile}}}
\fancyfoot[C]{\small\thepage}
\renewcommand{\headrulewidth}{0.8pt}

\begin{document}

\begin{center}
{\LARGE\bfseries\color{primarydark} Composées à une couche — corrigé}\\[2pt]
{\color{primary}\rule{0.5\textwidth}{1.2pt}}\\[4pt]
{\small\itshape Niveau facile}
\end{center}
\vspace{2mm}

\begin{corr}[1 --- $y=(2x+1)^5$]
Externe : puissance $5$. Intérieur $U=2x+1$, donc $U'=2$. Formule $(U^n)'=nU^{n-1}U'$ :
\[
y'=5(2x+1)^4\cdot 2 = 10(2x+1)^4.
\]
\end{corr}

\begin{corr}[2 --- $y=\sin(3x)$ puis $y=\cos(x^2)$]
$\bullet$ $U=3x$, $U'=3$, formule $(\sin U)'=U'\cos U$ :
\[
y'=3\cos(3x).
\]
$\bullet$ $U=x^2$, $U'=2x$, formule $(\cos U)'=-U'\sin U$ :
\[
y'=-2x\sin(x^2).
\]
\end{corr}

\begin{corr}[3 --- $y=e^{2x}$ puis $y=e^{-x^2}$]
$\bullet$ $U=2x$, $U'=2$, formule $(e^U)'=U'e^U$ :
\[
y'=2e^{2x}.
\]
$\bullet$ $U=-x^2$, $U'=-2x$ :
\[
y'=-2x\,e^{-x^2}.
\]
\end{corr}

\begin{corr}[4 --- $y=\ln(x^2+1)$]
$U=x^2+1$, $U'=2x$, formule $(\ln U)'=\dfrac{U'}{U}$ :
\[
y'=\frac{2x}{x^2+1}.
\]
(Défini sur $\R$ car $x^2+1>0$.)
\end{corr}

\begin{corr}[5 --- $y=\sqrt{x^2+1}$]
$U=x^2+1$, $U'=2x$, formule $(\sqrt U)'=\dfrac{U'}{2\sqrt U}$ :
\[
y'=\frac{2x}{2\sqrt{x^2+1}}=\frac{x}{\sqrt{x^2+1}}.
\]
\end{corr}

\end{document}
